\documentclass[letterpaper,11pt]{article}

\usepackage{latexsym}
\usepackage[empty]{fullpage}
\usepackage{titlesec}
\usepackage{marvosym}
\usepackage[usenames,dvipsnames]{color}
\usepackage{verbatim}
\usepackage{enumitem}
\usepackage[hidelinks]{hyperref}
\usepackage{fancyhdr}
\usepackage[english]{babel}
\usepackage{tabularx}
\input{glyphtounicode}

\pagestyle{fancy}
\fancyhf{} % clear all header and footer fields
\fancyfoot{}
\renewcommand{\headrulewidth}{0pt}
\renewcommand{\footrulewidth}{0pt}

% Adjust margins
\addtolength{\oddsidemargin}{-0.5in}
\addtolength{\evensidemargin}{-0.5in}
\addtolength{\textwidth}{1in}
\addtolength{\topmargin}{-.5in}
\addtolength{\textheight}{1.0in}

\urlstyle{same}

\raggedbottom
\raggedright
\setlength{\tabcolsep}{0in}

% Ensure that generate pdf is machine readable/ATS parsable
\pdfgentounicode=1

% Paragraph spacing and indentation
\setlength{\parindent}{0pt}
\setlength{\parskip}{0.5em}

\begin{document}

%----------HEADING----------
\begin{center}
    \textbf{\Huge \scshape Simon Chen} \\ \vspace{1pt}
    \small 347-322-4133 $|$ \href{mailto:shangminch@gmail.com}{\underline{shangminch@gmail.com}} $|$ \href{https://simon-chen.com}{\underline{simon-chen.com}} $|$ \href{https://linkedin.com/in/shangmin-chen}{\underline{linkedin.com/in/shangmin-chen}}
\end{center}

\vspace{10pt}

July 7, 2026

\vspace{8pt}
\textbf{Junior AI} \\
Talent Acquisition Team \\
New York, NY

\vspace{8pt}
\textbf{Re: Application for Full Stack Engineer}

\vspace{8pt}
Dear Annie Ntouva and the Junior AI Selection Committee,

Your note about the Full Stack Engineer role immediately caught my eye. I love building products from the ground up, so your focus on shipping LLM-powered software for private equity and consulting firms hits close to home. I work heavily with Next.js, TypeScript, and PostgreSQL. Because I enjoy owning features end-to-end, I would love to bring my background in data pipelines and performance optimization to your tight-knit team in NYC.

I look at software engineering as a balance between user experience and deep engineering stability. While working with Ezesports, I migrated our platform over to Next.js and a Supabase backend. By configuring edge caching and smart code splitting, I cut page load times by 75 percent. At RESET Standard, I handled background processing jobs by scaling out data pipelines with RabbitMQ and Redis. I know what it takes to move fast, ship clean code, and build infrastructure that businesses can actually trust.

A few independent projects show how I approach data and optimization:
\begin{itemize}[leftmargin=*,noitemsep,topsep=4pt]
    \item \textbf{LLMs and Media Pipelines:} I built \textit{Whisperrr}, an audio transcription platform that processes files four times faster than the stock OpenAI Whisper API by leveraging connection pooling and isolated Docker containers.
    \item \textbf{High-Throughput Backends:} I engineered \textit{Persephone}, a trading system that managed live order books. I pushed critical bottlenecks into native Cython extensions to drive latency down to the microsecond level.
    \item \textbf{Product Ownership:} I design responsive interfaces using React, Next.js, and Tailwind CSS. My goal is always to make complex data tools feel smooth and completely intuitive for the person using them.
\end{itemize}

Building real AI products with a lean, highly capable team is exactly the environment where I do my best work. I am based right here in New York and would love to stop by your office to chat about how my full-stack background can help you scale. Thanks for reaching out and for your time.

\vspace{10pt}
Sincerely,

\vspace{20pt}
Simon Chen

\end{document}
